\subsubsection{2.9. Nền tảng công nghệ sử dụng}

\indent\indent Để xây dựng một hệ thống tích hợp mô hình trí tuệ đồ thị với giao diện phân tích tương tác cao, nền tảng ứng dụng được chia thành hai phần chính: frontend phát triển trên framework Next.js và backend sử dụng FastAPI kết hợp hạ tầng điện toán đám mây Modal.\vspace{0.3cm}

Next.js\footnote{\url{https://nextjs.org/}} là một React framework mã nguồn mở được phát triển bởi Vercel, cung cấp các cơ chế kết xuất phía máy chủ (Server-Side Rendering -- SSR) và tạo trang tĩnh (Static Site Generation -- SSG). Việc sử dụng Next.js giúp tối ưu hóa tốc độ tải trang và duy trì sự mượt mà cho các bảng điều khiển phân tích dữ liệu phức tạp.\vspace{0.3cm}

Một thành phần công nghệ mang tính cốt lõi ở phía frontend là thư viện trực quan hóa \texttt{react-force-graph-2d}\footnote{\url{https://github.com/vasturiano/react-force-graph}}. Thư viện này sử dụng công cụ mô phỏng vật lý kết hợp với HTML5 Canvas để kết xuất cấu trúc mạng lưới chứa hàng ngàn nút và cạnh. Trong bối cảnh phát hiện Sybil, việc trực quan hóa giúp chuyển đổi các vector nhúng và ma trận dự đoán khô khan thành một bản đồ tri thức dễ tiếp cận. Các tính năng tương tác thời gian thực hỗ trợ người quản trị dễ dàng khoanh vùng các cụm tài khoản khả nghi hoặc mạng lưới đồng sở hữu, từ đó hiện thực hóa mục tiêu "khả năng diễn giải" (\textit{Explainability}) của toàn bộ đề tài.\vspace{0.3cm}

Về phần backend, FastAPI\footnote{\url{https://fastapi.tiangolo.com/}} được lựa chọn nhờ hiệu năng vượt trội và sự tương thích hoàn hảo với hệ sinh thái Python — ngôn ngữ của các thư viện học sâu và xử lý đồ thị. Hệ thống sử dụng Uvicorn làm ASGI server để tối ưu hóa khả năng xử lý bất đồng bộ, kết hợp cùng Pydantic nhằm đảm bảo tính chặt chẽ trong việc xác thực dữ liệu và quản lý cấu hình.\vspace{0.3cm}

Đối với luồng xử lý phân tích và AI, kiến trúc kết hợp NetworkX\footnote{\url{https://networkx.org/}} để thao tác trên các cấu trúc đồ thị phức tạp. Google Cloud BigQuery được tích hợp nhằm đáp ứng năng lực truy vấn với các tập dữ liệu quy mô lớn từ cơ sở dữ liệu của nền tảng Lens Protocol.\vspace{0.3cm}

Trong nền tảng backend, tích hợp hạ tầng điện toán phi máy chủ Modal\footnote{\url{https://modal.com/}} để triển khai mô hình và hệ thống các dịch vụ API, kết hợp cùng cơ sở dữ liệu SQLite để lưu vết lịch sử. Tổ hợp công nghệ này mang lại các ưu điểm vượt trội bao gồm:\vspace{-0.1cm}
\begin{itemize}[label={--}, itemsep=1pt]
    \item \textbf{Tối ưu hóa tính toán GPU theo nhu cầu:} Quá trình suy diễn với mạng nơ-ron đồ thị đa tầng đòi hỏi năng lực xử lý song song mạnh mẽ. Modal cung cấp môi trường điện toán đám mây cho phép khởi động lạnh (\textit{cold start}) nhanh chóng trên kiến trúc GPU và tính phí theo thời gian thực thi, giải quyết triệt để bài toán lãng phí tài nguyên khi phải duy trì máy chủ AI liên tục 24/7.
    \item \textbf{Xử lý bất đồng bộ hiệu năng cao:} Được xây dựng trên chuẩn ASGI (Asynchronous Server Gateway Interface), FastAPI cho phép xử lý hàng loạt yêu cầu phân tích đồ thị đồng thời mà không gây tắc nghẽn luồng (\textit{non-blocking}), giúp giảm thiểu tối đa độ trễ phản hồi khi tích hợp với các mô hình máy học.
    \item \textbf{Chuẩn hóa dữ liệu và tự động hóa tài liệu:} Hệ thống sử dụng Pydantic để định nghĩa và kiểm soát nghiêm ngặt cấu trúc dữ liệu đồ thị ở đầu vào và đầu ra, đồng thời tự động sinh tài liệu chuẩn OpenAPI (Swagger UI). Điều này đảm bảo tính nhất quán cao độ và hỗ trợ việc tích hợp trơn tru giữa ứng dụng client và các mô hình trí tuệ nhân tạo.
\end{itemize}